
\section{Chirality, antimatter, and \texorpdfstring{$CPT$}{CPT} as protocol geometry}
\label{sec:chirality_antimatter}

\subsection{Protocol-level \texorpdfstring{$P,T,C$}{P,T,C}: definitions}
\label{subsec:ptc_definitions}

\begin{definition}[Protocol-level discrete symmetries]\label{def:ptc_protocol}
Fix a Hilbert order $n$ and a window length $m$.
We define:
\begin{itemize}
\item $P_{\mathrm{prot}}$: a spatial reflection (an orientation-reversing element of the $D_4$ layout family acting on the Hilbert-addressed grid);
\item $T_{\mathrm{prot}}$: scan traversal reversal (replacing the scan index order $k\mapsto 4^n-1-k$ on the Hilbert path);
\item $C_{\mathrm{prot}}$: phase conjugation on the scan orbit, $z_n\mapsto \overline{z_n}$, with window conjugation chosen so that admissible readout is compared within matched protocols.
\end{itemize}
These are protocol transformations; they need not coincide with the continuum-field definitions of $P,T,C$ in standard QFT, but are designed to be auditable at finite resolution.
\end{definition}

\subsection{Scan--chirality locking}
\label{subsec:scl}

Hilbert addressing admits both orientation-preserving and orientation-reversing layout families (Section~\ref{sec:hilbert_addressing}).
We propose that a physical ``parity direction'' is fixed by the initialization of the readout protocol:

\begin{axiom}[Scan--chirality locking (SCL)]
\label{ax:scl}
At $t=0$, the readout protocol fixes an orientation class of Hilbert addressing.
Parity transformation (spatial reflection) therefore corresponds to a protocol change rather than a symmetry within a fixed protocol.
\end{axiom}

In this view, maximal parity violation in the weak sector is an effective statement: the weak compensation connection is defined only within the chosen protocol class, while its mirror is not available as an internal symmetry operation (it changes the readout basis).
\paragraph{Why this targets the weak sector.}
Within the three-channel template, the $\pi$-channel is the closure/topology channel, and the weak sector is the minimal non-abelian gauge sector $SU(2)$.
If the corresponding compensation connection is implemented by protocol-local consistency rules that depend on Hilbert layout orientation, then the mirror layout is not a symmetry but a different protocol.
This provides a readout-geometric route to maximal parity violation: the ``right-handed'' copy is not a dynamical possibility within the fixed protocol.

\subsection{Time reversal and discrete chirality}
\label{subsec:time_reversal_chi}

The chirality index $\chi$ in \eqref{eq:hilbert_chi_def} flips sign under traversal reversal (Proposition~\ref{prop:chi_flip}).
Since traversal reversal is the discrete operational avatar of time reversal in a scan process, this provides an auditable finite-resolution coupling:
\[
T\ \text{reversal} \quad\Longleftrightarrow\quad \chi\mapsto -\chi.
\]
Appendix~\ref{app:generated_tables} records the $n=3$ computation that $\chi(\text{path})=-2$ and $\chi(\text{reversed path})=+2$.

\paragraph{A discrete proxy for geometric phase.}
The index $\chi$ counts the signed turning of a locality-preserving address curve.
In continuous settings, signed turning and orientation couple naturally to geometric phase (Berry phase) and to spinorial sign structures \cite{Berry1984}.
The finite $n=3$ computation therefore provides a minimal auditable finite-resolution model for the statement
\[
T\ \text{reversal}\ \Longleftrightarrow\ \text{chirality/spin sign flip},
\]
at the protocol layer: both traversal reversal (time reversal) and reflection (parity) flip the same discrete sign datum.

\subsection{Antimatter as conjugate readout}
\label{subsec:antimatter_duality}

Consider the scan orbit $z_n=\e^{2\pi\iu(x_0+n\alpha)}$.
Complex conjugation yields
\[
\overline{z_n}=\e^{-2\pi\iu(x_0+n\alpha)}=\e^{2\pi\iu((-x_0)+(-n)\alpha)}.
\]
Thus phase conjugation is equivalent to scan reversal $n\mapsto -n$ up to a phase shift.
Formally:
\begin{lemma}[Conjugation equals reversal up to an initial-phase flip]\label{lem:conjugation_reversal}
Let $z_n=\e^{2\pi\iu(x_0+n\alpha)}$ and define $z'_n:=\e^{2\pi\iu(x_0'+n\alpha)}$ with $x_0'=-x_0$.
Then $\overline{z_n}=z'_{-n}$ for all $n\in\ZZ$.
\end{lemma}
\begin{proof}
This is the displayed computation.
\end{proof}
Consequently, in the rotation model the protocol operation $C_{\mathrm{prot}}$ is realized by traversal reversal together with an initial-phase flip, hence it is tied to the same discrete chirality sign datum as $T_{\mathrm{prot}}$:
\[
C_{\mathrm{prot}}\ \Longrightarrow\ \chi\mapsto -\chi.
\]
At the level of stable types, this motivates the definition of an antimatter dual $w\mapsto \overline{w}$ as the type induced by conjugate readout under matched windows.

\paragraph{Massless limit heuristic.}
In the relativistic massless limit, chirality coincides with helicity.
Under conjugation-as-reversal, a left-handed particle maps to a right-handed antiparticle, matching the operational observation that ``handedness'' is tied to traversal direction and protocol orientation.
\emph{Remark:} we use this only as a consistency heuristic with standard QFT kinematics; see, e.g., \cite{PeskinSchroeder1995}.

\subsection{\texorpdfstring{$CPT$}{CPT} at the scan layer vs.\ symmetry breaking at the protocol layer}
\label{subsec:cpt_protocol}

The scan layer is unitary and reversible by construction, so combined transformations can remain valid at the microscopic level.
However, once a readout protocol is fixed (window choice plus Hilbert orientation class), $P$ and $T$ may no longer be symmetries within the same protocol: they change the readout basis.
This provides a protocol-geometric route to understanding why effective theories may violate $P$ (and $CP$) while preserving $CPT$.

\subsection{Mirror protocols and a ``right-handed'' universe}
\label{subsec:mirror_universe}

If the initialization at $t=0$ selects the opposite Hilbert orientation class, then the protocol-level chirality sign flips globally.
In such a mirror protocol, the weak sector would appear right-handed rather than left-handed.
This provides a precise sense in which a ``mirror universe'' is not a different Lagrangian, but a different readout protocol (Hilbert layout class).
If protocol domains existed in the early universe, their boundaries would be protocol defects, motivating the domain-wall prediction in Section~\ref{subsec:p2_domain_walls}.


